% Theoretical Framework
% =============================
% This section reviews the main concepts: continuous assessment, fuzzy logic in education, and student-centered learning.
% TODO: Expand with more recent studies and add figures if needed.

\section{Theoretical Framework}

\subsection{Continuous Assessment in Education}
Continuous assessment is a formative approach that emphasizes ongoing feedback and the integration of multiple sources of evidence to support student learning. Recent studies highlight its effectiveness in promoting engagement and deeper understanding, especially in remote and hybrid environments \cite{Lopez2021, Black2020}.

\subsection{Fuzzy Logic in Educational Assessment}
Fuzzy logic provides a mathematical framework to handle uncertainty and subjectivity in educational evaluation. Its application allows for the integration of qualitative and quantitative variables, making grading fairer and more transparent. Recent research demonstrates the potential of fuzzy systems to automate and personalize assessment processes \cite{Santos2022, Wang2023}.

\subsection{Student-Centered Learning}
Student-centered learning shifts the focus from teaching to learning, prioritizing the development of competencies, autonomy, and critical thinking. This paradigm is especially relevant in digital education, where flexibility and adaptability are crucial \cite{Nguyen2020, Smith2023}.

% TODO: Add a conceptual diagram of the fuzzy assessment system in /figuras

The integration of continuous assessment, fuzzy logic, and student-centered learning forms the basis for innovative and equitable evaluation models in contemporary education.
